\documentclass[DIV=calc, paper=a4, fontsize=11pt, twocolumn]{scrartcl}	 % A4 paper and 11pt font size

\usepackage{lipsum} % Used for inserting dummy 'Lorem ipsum' text into the template
\usepackage[english]{babel} % English language/hyphenation
\usepackage[protrusion=true,expansion=true]{microtype} % Better typography
\usepackage{amsmath,amsfonts,amsthm} % Math packages
\usepackage[svgnames]{xcolor} % Enabling colors by their 'svgnames'
\usepackage[hang, small,labelfont=bf,up,textfont=it,up]{caption} % Custom captions under/above floats in tables or figures
\usepackage{booktabs} % Horizontal rules in tables
\usepackage{fix-cm}	 % Custom font sizes - used for the initial letter in the document
\usepackage{graphicx} % Package to insert images

\usepackage{sectsty} % Enables custom section titles
\allsectionsfont{\usefont{OT1}{phv}{b}{n}} % Change the font of all section commands

\usepackage{fancyhdr} % Needed to define custom headers/footers
\pagestyle{fancy} % Enables the custom headers/footers
\usepackage{lastpage} % Used to determine the number of pages in the document (for "Page X of Total")

% Headers - all currently empty
\lhead{Physics 11422}
\chead{}
\rhead{Spring 2020}

% Footers
%\lfoot{\footnotesize Instructor: } % Instructor's Name
%\cfoot{\footnotesize TA: } % TA's Name
\rfoot{\footnotesize Page \thepage\ of \pageref{LastPage}} % Format for Footnote: Page 1 of 2

\renewcommand{\headrulewidth}{0.0pt} % No header rule
\renewcommand{\footrulewidth}{0.4pt} % Thin footer rule

\usepackage{lettrine} % Package to accentuate the first letter of the text
\newcommand{\initial}[1]{ % Defines the command and style for the first letter
\lettrine[lines=3,lhang=0.3,nindent=0em]{
\color{DarkGoldenrod}
{\textsf{#1}}}{}}


%----------------------------------------------------------------------------------------
%	TITLE SECTION
%----------------------------------------------------------------------------------------

\usepackage{titling} % Allows custom title configuration

\newcommand{\HorRule}{\color{DarkGoldenrod} \rule{\linewidth}{1pt}} % Defines the gold horizontal rule around the title

\pretitle{\vspace{-30pt} \begin{flushleft} \HorRule \fontsize{50}{50} \usefont{OT1}{phv}{b}{n} \color{DarkRed} \selectfont} % Horizontal rule before the title

% Lab Report Title Goes Here
\title{Lab Report Title} 
% Lab Report Title Goes Here

\posttitle{\par\end{flushleft}\vskip 0.5em} % Whitespace under the title

\preauthor{\begin{flushleft}\large \lineskip 0.5em \usefont{OT1}{phv}{m}{sl} \color{Black}} % Author font configuration

% Student's name(s)
\author{Student 1, Student 2, Student 3, } 
% Student's name(s)

% Student's Institution
\postauthor{\footnotesize \lineskip 0.5em \usefont{OT1}{phv}{m}{sl} \color{Black} % Configuration for the institution name
University of Notre Dame\\ 
% Student's institution

\par\end{flushleft}\HorRule} % Horizontal rule after the title

\date{} % Add a date here if you would like one to appear underneath the title block

%----------------------------------------------------------------------------------------

\begin{document}

\maketitle % Print the title

\thispagestyle{empty} % Enabling the custom headers/footers for the second page forward

%----------------------------------------------------------------------------------------
%	ABSTRACT
%----------------------------------------------------------------------------------------

% The first character should be within \initial{}
\initial{T}\textbf{is the space where students place their abstract. Please see the How to Write a Lab Report Document and the Lab Report Rubric. In the lab report guidelines, students will find MatLib Abstract.}

%----------------------------------------------------------------------------------------
%	ARTICLE CONTENTS
%----------------------------------------------------------------------------------------

\section*{Introduction/Brief Theory Review}

This Introductory Physics laboratory course has been developed with two goals in mind. First, the Physics lab should help students to understand the concepts presented in lecture, and to reinforce this understanding by experience. Second, since Physics is an experimental science, the lab should introduce the students to the methodology of planning and executing an experiment. The students should be able to analyze the situation, to provide a coherent writeup of the purpose and manner in which they understood the lab, and to estimate the uncertainties/errors within their process, measurements, and analysis.

This first section should include a brief description of the physics knowledge needed to understand this Experimental Design (typically including some basic equations). Here is a mathematical equation that will automatically number equations and matrices. Edit or delete this for your own lab report. Please see LaTex formula or cheat sheets as they are called.

\begin{align}
\cos (2\theta) = \cos^2 \theta - \sin^2 \theta
\end{align}

Here is another type of equation with subscripts.

\begin{align}
k_{n+1} = n^2 + k_n^2 - k_{n-1}
\end{align}

Here is a matrix that will automatically number equations and matrices.  Edit or delete this for your own lab report.

\begin{align}
A = 
\begin{bmatrix}
A_{11} & A_{21} \\
A_{21} & A_{22}
\end{bmatrix}
\end{align}

\begin{figure}
  \includegraphics[width=\linewidth]{physicsimage.jpeg}
  \caption{Physics knowledge encompasses the extremely large (relativity) to small (quantum mechanics)}
  \label{fig}
\end{figure}

Sometimes figures are good to use to explain a physics concept as seen in Figure 1.

%------------------------------------------------

\section*{Experimental Design}

Make sure to clearly explain how you designed the laboratory in terms of what is measured and why it was measured. Do not include the actual measurements in this section, but rather in the the Experimental Observation/Data section. This section is only to explain enough details for other physicists to repeat or slightly change the procedures to execute the experimental design (again, do not include the data here... that comes later on).

%------------------------------------------------

\subsection*{Apparatus}

There should always be an image in the Apparatus section. Below is the LaTex code, but remember that in LaTex, this code will format the image placing in the document (in this case top left corner of document). In this template, the sample apparatus figure is labeled Figure 2.

\begin{figure}
  \includegraphics[width=\linewidth]{apparatuspicture.jpg}
  \caption{Please look at Google Images for formal physics apparatus images, but MAKE SURE TO REFERENCE THESE. Edit the images so that Figure 1 is not in the image.}
  \label{fig2}
\end{figure}

This Apparatus subsection should describe VERY BRIEFLY the apparatus (2-3 sentences). An appendix can be referenced and added to describe unique or novel apparatus discussions.

\subsection*{Design Matrix/Trials}

This is typically used to highlight a basic understanding of an experimental design in terms of a design matrix (what was varied) and what was not varied (so called controlled or assumed to be controlled variables/conditions of the experiments).

\begin{table}
\caption{Title of Design Matrix}
\centering
\begin{tabular}{llr}
\toprule
\multicolumn{2}{c}{Design Matrix} \\
\cmidrule(r){1-2}
Variable & Units & Trials \\
\midrule
Position & meters & $10$ \\
Time & seconds & $10$ \\
Incline & degrees & $6$ \\
\bottomrule
\end{tabular}
\end{table}


\lipsum[3] % Dummy text

%------------------------------------------------

\section*{Experimental Observations/Data}

This section will only include the raw data measured. Consequently no calculations should be shown. Many times a chart or graph is used to present the data and then the full dataset is presented in an Appendix.

%------------------------------------------------

\section*{Experimental Analysis}

\lipsum[3] % Dummy text

This equation has a fraction within a fraction.

\begin{align}
\psi_{1}(x,y,z) = \frac{\frac{1}{x}+\frac{1}{y}}{y-z}
\end{align}

Here is one that looks the same, but LaTex uses begin equation instead of begin align in the LaTex code. Remember all equations should be derived from first principles. If extra space is needed, please reference an Appendix.

\begin{equation} 
x = a_0 + \cfrac{1}{a_1 + \cfrac{1}{a_2 + \cfrac{1}{a_3 + \cfrac{1}{a_4}} } } 
\end{equation}

\begin{figure}
  \includegraphics[width=\linewidth]{physicsimage.jpg}
  \caption{Make sure to label all graphs/charts correctly before inserting into LaTex Lab Report. This is a sample fun results image.}
  \label{fig3}
\end{figure}

This analysis section will also describe the three uncertainties/errors after a well-crafted results graph/chart, see Figure 3. The three types of uncertainties/errors are process, measurement, and analysis. They can be either random or systematic and a brief discussion and final uncertainty/error calculation should be written up. The abstract should have a result sentence, see MadLib for format of this type of sentence, with uncertainty/error included.

%------------------------------------------------

\section*{Conclusion}

Conclusion concisely summarizes the laboratory (experimental design, experimental observation, experimental analysis). It evaluates the results relative to “what the goal/hypothesis of the experiment was” and relative to the uncertainties/errors.  A quick discussion of how one could make the laboratory better to reduce uncertainties(random)/errors (systematic) is vital. This should be 3-7 sentences.


%----------------------------------------------------------------------------------------
%	REFERENCE LIST
%----------------------------------------------------------------------------------------

\begin{thebibliography}{10} % 10 is a random guess of the total number of references

\bibitem{MG} Goossens, M., Mittelbach, F., Samarin, \emph{A LaTeX
Companion}, Addison-Wesley, Reading, MA, 1994.

\bibitem{HK} Kopka, H., Daly P.W., \emph{A Guide to LaTeX},
Addison-Wesley, Reading, MA, 1999.

\bibitem{PR} Paul Reimer, \emph{http://www.phy.anl.gov/mep/SeaQuest/}, 2014. 

\bibitem{JA} Joseph C. Amato and Roger E. Williams, \emph{AAPT Apparatus Competition}, 2008. 

\end{thebibliography}

%----------------------------------------------------------------------------------------

\newpage 

%----------------------------------------------------------------------------------------
%	APPENDIX CONTEXT
%----------------------------------------------------------------------------------------

\section*{Appendix A}

Add extra information here or delete this section.

\end{document}
